\section{Líneas de trabajo futuro}
\label{sec:lineas_futuras}

El desarrollo de Voctomix~2.0 ha puesto de manifiesto varias áreas de mejora e
investigación que quedan fuera del alcance de este Trabajo de Fin de Grado pero que
constituyen líneas de trabajo futuro de interés.

\begin{itemize}

  \item \textbf{Validación con cámaras físicas en producción real.}
  La totalidad de las pruebas descritas en este trabajo se han realizado con fuentes
  sintéticas generadas mediante FFmpeg. La validación definitiva del sistema consiste
  en conectar cámaras físicas mediante tarjetas de captura o NDI y realizar una
  retransmisión en directo completa. Este escenario está previsto en el contexto del
  grupo de investigación CyberNemo de la UPM, donde Voctomix~2.0 se utilizará para
  retransmitir seminarios y presentaciones. Una producción real permitirá detectar
  posibles problemas de sincronización, \textit{jitter} de red o estabilidad que no
  se manifiestan con fuentes sintéticas de frecuencia constante.

  \item \textbf{Migración del pipeline de salida de H.264 a H.265/HEVC.}
  Voctomix~2.0 emplea H.264 como códec de salida predeterminado a través de
  \texttt{libx264}, que opera en tiempo real en CPU de gama media. La migración del
  pipeline de transcodificación a H.265/HEVC permitiría reducir el ancho de banda de
  la señal de programa aproximadamente a la mitad manteniendo la misma calidad
  perceptual, lo que facilitaría la distribución en entornos con conectividad
  limitada. El principal obstáculo es la velocidad de codificación, entre cinco y
  ocho veces inferior a la de H.264, por lo que la implementación en tiempo real
  requeriría aceleración hardware mediante NVENC (NVIDIA) o VAAPI (Intel), ambas
  accesibles desde FFmpeg. Esta mejora es directamente aplicable en los escenarios
  contenerizados, donde el contenedor de codificación puede actualizarse de forma
  independiente sin modificar el núcleo de voctocore.

  \item \textbf{Escalado de la resolución de trabajo a 4K.}
  La resolución de trabajo actual del sistema es 1920$\times$1080 píxeles a 25~fps,
  con un ancho de banda interno de 622~Mbit/s por fuente RAW~I420. La ampliación a
  resolución 4K (3840$\times$2160) cuadruplicaría ese ancho de banda hasta
  aproximadamente 2{,}5~Gbit/s por fuente, lo que exigiría circunscribir el tráfico
  de vídeo a redes de alta velocidad o interfaces de red con MTU ampliado. El motor
  compositor de GStreamer y los modos de composición existentes admiten resoluciones
  arbitrarias, de forma que la adaptación se realizaría principalmente mediante
  cambios en el fichero de configuración de voctocore, si bien los requisitos de CPU
  aumentarían proporcionalmente al número de píxeles procesados. Este escalado
  situaría el sistema a la par con los estándares de producción actualmente exigidos
  en eventos de mediana y gran escala.

  \item \textbf{Interfaz gráfica con reconfiguración dinámica en tiempo real.}
  La arquitectura actual de voctogui requiere reiniciar el sistema completo para
  modificar el número de fuentes activas o la configuración de los modos de
  composición, lo que limita la adaptabilidad del sistema durante una producción en
  directo. Una extensión natural sería ampliar el protocolo de control TCP del
  puerto~9999 con comandos que permitan añadir y eliminar fuentes en caliente,
  modificar los parámetros de los composites en tiempo real y recargar la
  configuración sin interrumpir la señal de programa. En el lado del cliente, la
  interfaz gráfica debería reflejar dinámicamente esos cambios actualizando sus
  controles sin necesidad de reconexión. Esta mejora eliminaría uno de los
  principales puntos de fricción del sistema en producciones con configuraciones
  variables o imprevistos técnicos durante la emisión.

\end{itemize}
